\documentclass[9pt,twocolumn]{article}
\usepackage[letterpaper,margin=0.55in,top=0.48in,bottom=0.52in]{geometry}
\usepackage[T1]{fontenc}\usepackage{newtxtext,newtxmath}\usepackage{microtype}
\usepackage{amsmath,mathtools}\usepackage{booktabs,tabularx,array}\usepackage[normalem]{ulem}\usepackage{cancel}\usepackage{enumitem,xcolor,ragged2e}
\usepackage{graphicx}
\setlength{\columnsep}{0.24in}\setlength{\parindent}{0pt}\setlength{\parskip}{2.2pt}
\setlist{nosep,leftmargin=*}\renewcommand{\arraystretch}{1.10}
\newcommand{\oldclaim}[1]{\sout{#1}}
\newcommand{\claimbox}[1]{\begin{center}\fbox{\parbox{0.94\columnwidth}{\centering\scriptsize #1}}\end{center}}
\newcommand{\eclass}[1]{\textsf{[#1]}}
\DeclareFontFamilySubstitution{TS1}{ntxtlf}{cmr}
% \let binds the CURRENT meaning of \_ ; \newcommand would create a macro that
% re-expands to \_ and loops until TeX runs out of input stack.
\let\CIGunderscore\_
\renewcommand{\_}{\CIGunderscore\allowbreak}
\begin{document}
\sloppy
\twocolumn[{\begin{center}
{\LARGE\bfseries Report 1 - MASTER Cross-Chain Synthesis}\\[2pt]
{\large\itshape Append-only update D: three verifications, one proposed target shown weaker than the evidence in hand, and a corpus-wide trust-hole census}\\[2pt]
{\small CIG-AMF Scientific Evolution Report -- 8 September (third update)}
\end{center}\vspace{2mm}}]
\textbf{Append-only rule.} Sections 1--67, including addenda A--C, are an immutable prefix. This addendum begins at Section 68.

\section*{68. What was verified this cycle}
Three checks, each stated with its evidence class rather than its conclusion.

\textbf{(a) Block composition.} For $F(x)=\sum_b F_b(x_b)$ with the product reference, $\delta_\square(F)=\max_b\delta_\square(F_b)$ and $2D(F)=\sum_b 2D(F_b)$. Composing two copies of the exact $m=4$ witness gave $\delta_\square=1$ and $2D=6$, matching both predictions exactly over $\mathbb{Q}$. \eclass{VERIFIED NUMERICALLY, unformalised} --- see D6 Section 147.

\textbf{(b) The information limit has a closed form.} On the family generated by permuting the $m-1$ positive-slope coordinates --- worlds with \emph{identical} response profile $Q$ and different optima --- the game value is exactly
\[V=k-\frac{k^2}{m-1}=\frac{k(k-1)}{2k-1},\]
confirmed by exact LP at $m=6,8,10,12$ (values $6/5,\,12/7,\,20/9,\,30/11$). \eclass{EXACT LP on a constructed family; a lower bound for the class, not a minimax solution} --- see D6 Section 148.

\textbf{(c) Trust-hole census.} Nine theorems in three modules are proved by \texttt{native\_decide}, a tactic this project rejects. \eclass{COMPILED ARTIFACT, not PROVED} for those nine --- see Structural Section 67.

\section*{69. NARROWING a target proposed for the D6 information limit}
\begin{center}\scriptsize\oldclaim{nearest target: $V\ge m/5$ at $\mu=0$ for $m=6r$, via six-dimensional blocks plus a global budget argument}\end{center}
$k(k-1)/(2k-1)=m/5$ only at $m=6$; at $m=8$ the family already gives $12/7\approx1.714$ against $1.6$, and the closed form tends to $m/4$. The family is moreover defined \emph{directly} at every even $m$, so this bound needs \textbf{no block composition and no $m=10$ witness}. A composition route would be both unnecessary and lossy, recovering $m/5$ where $m/4$ is available.

\claimbox{\textbf{The two D6 programmes must not be merged.} The Top-C frontier target $R_m(\mu)=m-\tfrac m2\mu$ \emph{does} require composition and \emph{does} have an $m=10$ gap, since $\langle6,8\rangle$ misses $10$. The information limit does not. Treating them as one programme imports the $m=10$ obstruction into a statement that never had it, and costs the constant.}

\section*{70. A method error of mine, recorded}
Asked how many modules use \texttt{native\_decide}, I first grepped for the string \texttt{by native\_decide} and reported one module. That is wrong: in two of the three modules the tactic sits on the line \emph{after} \texttt{by}, so the pattern misses them. I also initially counted docstring mentions as uses, including two docstrings that say the tactic is \emph{not} used --- an error in the opposite direction.
\claimbox{\textbf{Rule, added to the standing method notes.} Count occurrences of the bare identifier, then classify each occurrence by reading its context: real tactic, docstring, or comment. Never let a two-token grep decide an evidence class. This is the same failure mode as Section 52 (reading a result artifact without its companion status record) --- a cheap query substituted for the check that was actually needed.}

\section*{71. Where each chain now stands}
\begin{center}\scriptsize
\begin{tabularx}{0.98\columnwidth}{l X X}
\toprule
Chain & Best established & Immediate next step\\
\midrule
D6 & $2[L_F(S)-L_F(T)]\le d(2\delta_\square-\mu)$ \eclass{PROVED}; $c^\ast(4)=3$, $c^\ast(6)=6$, $c^\ast(8)=8$ exact; $R_6(\mu)=6-3\mu$ at five $\mu$; $V\ge k(k-1)/(2k-1)$ & formalise: permutation preserves $Q$ $\to$ regret $=$ distance to optimal set $\to$ lower bound for every algorithm\\
Structural & $\chi_{\mathrm{prefix}}(F_r,0)=r=m/2$ \eclass{PROVED} parametrically; staircase \textbf{not} extremal ($\chi\ge5$ at $m=7$) & replace $5\times$ \texttt{native\_decide} with kernel \texttt{decide}; then retarget to representability\\
Functional & both allocator defects real and repaired; $55$ anytime wins vs $0$ within-instance across $3$ seeds & $300$-instance confirmation at fixed budget\\
Support & fibre sufficiency: pairwise separation \emph{insufficient} \eclass{PROVED}, $10$ theorems audited & scope-check against the classical identification literature\\
Query & $h=1,4,8$ all give $0$ evaluable / $336$ degenerate; root cause is the competence of $\rho$, not the horizon & a continuation policy that is not constant NOOP\\
MASTER & $142$ modules / $1016$ theorems / $139$ compiled / $62$ audited; $0$ \texttt{sorryAx}, $0$ \texttt{ofReduceBool} among audited & audit the $9$ \texttt{native\_decide} theorems into the log so the hole is visible, not absent\\
\bottomrule
\end{tabularx}
\end{center}

\section*{72. Novelty ledger, carried forward}
Recorded so no cycle re-derives a known object: mixed differences and additive decomposition are classical; antichains of distinct-size sets sit in a known combinatorial circle; ``one ordering for every budget'' is studied in incremental submodular maximisation; variance-aware allocation and shrinkage estimation are standard. What remains defensible is narrow and should be stated narrowly: \emph{witnesses attaining a boundary}, the \emph{information limit} with its closed form, the \emph{insufficiency} of pairwise separation, and the \emph{representability} question for budget-indexed optima.
\end{document}
